\chapter{Superseded Permutations Draft}
\label{ch:permutations}

A vector register holds several words side by side, and the instructions that move those words between lanes --- shuffles, permutes, unpacks, blends --- are the ones this book's first machine theorems are about. This chapter sets up what a permutation of lanes is, why solving for one is a pure Boolean problem, and why the lower bounds that follow cannot come from counting.

\section{Lanes as tags}
\label{sec:tags}

The obvious model of a 256-bit register is thirty-two bytes. It is the wrong model for proving things about shuffles, for a reason that is easy to state and easy to forget: two bytes can be equal. If the register holds the value \(0\) in lanes three and seventeen, a shuffle that swaps those lanes is indistinguishable from one that does not, and a test that runs the shuffle on that register proves nothing about which permutation was applied.

\begin{artifact}{P.def.lane-tags}
Model the register instead as thirty-two \emph{distinct} symbolic tags, one per input lane. An instruction is a function on tag vectors. Replaying an instruction sequence on tags establishes the complete permutation --- and every zero the sequence produces --- in a single run, because no two tags can collide. A model that fails to replay is rejected outright; a mismatch between replay and the solver's model is a soundness alarm, not a corner case.
\ArtifactScope{P.def.lane-tags}
\end{artifact}

The consequence for proving minimality is the one that makes Part IV possible. Because tags are symbolic and the instructions are permutation-preserving, no bit-vector arithmetic enters the encoding at all. Each lane at each step is a one-hot choice among thirty-two tags; each instruction's control operand is a one-hot choice among its legal values; the constraints are that the selected instruction's semantics hold lane by lane. This is pure propositional logic. It needs no bit-blasting layer, and a proof-producing conflict-driven solver emits its refutation as DRAT natively.

\section{Why counting cannot give the lower bound}
\label{sec:entropy}

Here is the first place a reader might expect an easy argument and not find one. There are \(32!\) permutations of thirty-two lanes. A single \code{vpshufb} takes a thirty-two-byte control mask. Could an information-theoretic argument show that one instruction cannot reach every permutation?

\begin{artifact}{P.comp.entropy-bound}
No. \(\log_2 32! \approx 117.7\) bits, and a \code{vpshufb} control mask carries \(32 \times 4 = 128\) bits --- enough to name any permutation of thirty-two lanes with ten bits to spare. So the fact that byte reversal needs two instructions is not a counting fact. It is a \emph{lane-structure} fact: \code{vpshufb} cannot move a byte across the 128-bit lane boundary, and \code{vperm2i128} cannot move a byte within one. That is precisely the kind of fact that resists a pencil and yields to a solver.
\ArtifactScope{P.comp.entropy-bound}
\end{artifact}

\begin{keyidea}[title={Immediates versus registers}]
The same observation splits the instruction set into two kinds. Instructions whose control is an \emph{immediate} --- \code{vpermq}, \code{vperm2i128}, \code{vpalignr} --- have a small finite set of behaviours; the group they generate can be enumerated by breadth-first search, and its diameter computed exactly. Instructions whose control is a \emph{register} --- \code{vpshufb}, \code{vpermd}, \code{vpermb}, NEON's \code{tbl}, RISC-V's \code{vrgather} --- have astronomically many behaviours (\(16^{16}\) for a single \code{vpshufb} lane group), and breadth-first search is dead on arrival. For those, a SAT encoding with the control mask as free variables is the only formulation that closes. This is the cleanest justification available for why the theorems of Part IV needed a solver rather than a table.
\end{keyidea}

\section{Permutation networks}
\label{sec:networks}

The theory of permutation networks --- Beneš, Waksman --- says how many switching stages an arbitrary permutation of \(n\) elements needs: \(2\log_2 n - 1\), so nine for thirty-two lanes, with \(n\log_2 n - n + 1 = 129\) switches. RISC-V's Bitmanip draft uses exactly this to bound arbitrary bit permutations at eighteen instructions on a thirty-two-bit word. What the network literature does not supply is the transfer to a real instruction set: a stage is not an instruction, a switch is not a lane, and no theorem says how many \emph{instructions} a lane-restricted ISA needs for an arbitrary permutation. That transfer is an open problem in this book's catalogue, and the only route to a lower bound that scales without a solver run per permutation.

\section{Exercises}

\begin{enumerate}
\item Compute the group generated by same-source \code{vperm2i128} with all immediates, and its diameter. Is it a subgroup of \(S_{32}\)? Of what order?
\item Show that the entropy bound \emph{does} obstruct a single \code{vpermq}: how many bits does its immediate carry, and how many distinct permutations can it therefore express?
\item Find a permutation of thirty-two lanes that \emph{one} \code{vpshufb} can perform and \emph{one} \code{vpermd} cannot, and one for which the reverse holds. Replay both on tags.
\end{enumerate}
